\chapter[Band 34: Linkskürzbare Halbgruppen]{Band 34 auf einen Blick}
\noindent\textbf{Linkskürzbare Halbgruppen}\par\medskip
Vorausgesetzt sind Halbgruppen und die natürliche Addition.
Links- und Rechtskürzbarkeit werden als getrennte Eigenschaften geführt;
in Band 36 werden beide gemeinsam verlangt.

\begin{BandUebersicht}
\textbf{Kürzungsaxiomatik}
Zur Halbgruppenstruktur kommt die angegebene Kürzungsrichtung:\UebFormel{\begin{aligned}
&a,b,c\in A,\ a\star b=a\star c\ \vdash b=c.
\end{aligned}}Das Strukturprädikat lautet \(\LeftCancellativeSemigroupAlg A\star\).
\UebQuellen{\FormulaRefAuto{LeftCancellativeSemigroupDef}[def]}
&
\textbf{Die natürliche Addition}
\UebFormel{\begin{aligned}
&\LeftCancellativeSemigroupAlg{\mathbb N}{+},\\
&a,b,c\in\mathbb N,\ a+b=a+c\ \vdash b=c.
\end{aligned}}
\UebQuellen{\FormulaRefAuto{NaturalAdditionLeftCancellativeSemigroup}[thm];
\FormulaRefAuto{NaturalAdditionLeftCancellative}[thm]} \\
\addlinespace[.8ex]

\textbf{Homomorphismen und Isomorphismen}
Für zwei Strukturen derselben Klasse bedeutet ein Homomorphismus
\(H(f)\):\UebFormel{\begin{aligned}
&f:A\to B,\\
&x,y\in A\ \vdash\\
&\quad f(x\star y)=f(x)\diamond f(y).
\end{aligned}}Ein Isomorphismus \(I(f)\) ist ein bijektiver Homomorphismus:\UebFormel{\begin{aligned}
&I(f)\ \leftrightarrow H(f)\land f:A\bij B.
\end{aligned}}
\UebQuellen{\FormulaRefAuto{LeftCancellativeSemigroupHomDef}[def];
\FormulaRefAuto{LeftCancellativeSemigroupIsoDef}[def]}
&
\textbf{Komposition erhält die Struktur}
Für Homomorphismen \(f:A\to B\), \(g:B\to C\) zwischen
Strukturen dieser Klasse ist \(g\circ f\) wieder ein Homomorphismus.
Mit \(\bullet\) als Operation auf \(C\) gilt für \(x,y\in A\):\UebFormel{\begin{aligned}
&(g\circ f)(x\star y)\\
&\quad=(g\circ f)(x)\mathbin{\bullet}(g\circ f)(y).
\end{aligned}}Sind \(f,g\) Isomorphismen, so ist auch \(g\circ f\) ein Isomorphismus.
\UebQuellen{\FormulaRefAuto{SemigroupHomComposition}[thm];
\FormulaRefAuto{SemigroupIsoComposition}[thm]} \\
\addlinespace[.8ex]
\end{BandUebersicht}

